
\begin{Quizz}
    L'identifiant de variable \textit{1prenom} est-il valide ?
    \begin{itemize}
    \renewcommand{\labelitemi}{$\square$}
    \item Oui
    \item Non
    \end{itemize}
\end{Quizz}

\begin{Quizz}
    L'identifiant de variable \textit{\_prenom} est-il valide ?
    \begin{itemize}
    \renewcommand{\labelitemi}{$\square$}
    \item Oui
    \item Non
    \end{itemize}
\end{Quizz}

% \begin{Quizz}
%     Le nom de la variable suivante est-il valide ?
%     \textit{PreNom}
%     \begin{itemize}
%     \renewcommand{\labelitemi}{$\square$}
%     \item Oui
%     \item Non
%     \end{itemize}
% \end{Quizz}

\begin{Quizz}
     \begin{lstlisting}[language=Python]
a = 10     \end{lstlisting}
    Cochez les réponses vraies :
    \begin{itemize}
    \renewcommand{\labelitemi}{$\square$}
    \item a vaut 10
    \item a et 10 sont des variables égales
    \item La valeur de a est 10
    \item a est égale à 10
    \end{itemize}
\end{Quizz}


\begin{Quizz}
    \begin{lstlisting}[language=Python]
age = 15 
print("J'aurai", age+1, " ans l'année prochaine")\end{lstlisting}
    Quel est le type de la variable age ?
    \begin{itemize}
    \renewcommand{\labelitemi}{$\square$}
    \item int
    \item float
    \item str
    \item bool
    \end{itemize}
\end{Quizz}



\begin{Quizz}
    Pour calculer le poids en physique, on sait que $P=mg$. On écrit un script python où on demande à l'utilisateur sa masse en kg et on lui affiche son poids en Newton. Quels noms de variables sont les plus adaptées pour les variables de la formule ?
    \begin{itemize}
    \renewcommand{\labelitemi}{$\square$}
    \item poids, masse, G
    \item p,m,g
    \item Poids, Masse, g
    \item poids, masse, gravité
    \end{itemize}
\end{Quizz}


%Les opérateurs